\chapter{STM32 Pin Tahsisi ve Bağlantı Şeması}
\label{app:pin-plani}

\section{Pin Tahsis Tablosu (NUCLEO-F767ZI)}

\begin{table}[h!]
\centering
\caption{STM32 NUCLEO-F767ZI çevre birim ve GPIO tahsisleri (Plan B v2).}
\label{tab:pin-tahsisi}
\small
\begin{tabular}{l l l l}
\toprule
\textbf{İşlev} & \textbf{Pin} & \textbf{Bağlantı} & \textbf{Not}\\
\midrule
\multicolumn{4}{l}{\emph{USART6 --- ESP32-CAM iletişimi (115200 bps 8N1)}}\\
TX & PG14 (D1, Arduino başlık)  & ESP32-CAM RX (IO14) & Komut + telefon-yönlü mesaj\\
RX & PG9  (D0, Arduino başlık)  & ESP32-CAM TX (IO13) & RESULT/ERR/HB satırı\\
\midrule
\multicolumn{4}{l}{\emph{USART3 --- ST-LINK VCP debug (115200 bps 8N1)}}\\
TX & PD8 & ST-LINK (Mac terminal)       & \texttt{printf} retarget\\
RX & PD9 & (Kullanılmıyor)               & ---\\
\midrule
\multicolumn{4}{l}{\emph{Buton girişleri (EXTI)}}\\
TARA  & PC13 & B1 USER (kart üstü mavi buton)        & EXTI13, pull-up\\
PANİK & PA0  & Harici buton (sergi düzeneğinde takılı değil) & EXTI0, pull-up\\
\midrule
\multicolumn{4}{l}{\emph{GPIO çıkışları (push-pull)}}\\
LD1 yeşil   & PB0  & NUCLEO onboard            & IDLE / NOMATCH göstergesi\\
LD2 mavi    & PB7  & NUCLEO onboard            & Heartbeat (\~1 Hz) + NETERR blink\\
LD3 kırmızı & PB14 & NUCLEO onboard            & MATCH / PANIC\\
Buzzer      & PD14 & Aktif buzzer 5\,V (transistor sürücülü) & MATCH / PANIC\\
\midrule
\multicolumn{4}{l}{\emph{Kapalı / kullanılmayan çevre birimleri}}\\
ETH      & --- & (Disable)        & DCMI pin çakışması; Plan~B'de DCMI N/A\\
DCMI     & --- & (Disable)        & Kamera ESP32-CAM tarafında\\
USART1   & PA9/PA10 & (Boş)       & Morpho başlığı, kullanılmıyor\\
USART2   & PD5/PD6  & (Boş)       & HM-10 atıldığı için kullanılmıyor\\
\bottomrule
\end{tabular}
\end{table}

\section{Saat Üretimi}

Kart üzerindeki HSE 8~MHz osilatörü tercih edilmemiş; CubeMX'in ürettiği
\texttt{stm32f7xx\_hal\_conf.h} dosyasında \texttt{HSE\_VALUE} makrosu
25~MHz olarak tanımlandığı için (NUCLEO-F767ZI için yanlış varsayılan;
\cite{stm32f767_rm}) HSE üzerinden PLL yapıldığında HAL'in UART BRR
hesabı sapıyor, baud hızı yaklaşık 1/3 oranında bozuluyordu. Onun yerine dahili HSI 16~MHz
osilatörü PLL girişine bağlanır; PLL\_M=8, PLL\_N=216, PLL\_P=2,
PLL\_Q=9 yapılandırmasıyla SYSCLK $=$ HCLK $=$ 216~MHz elde edilir.
APB1 = 54~MHz, APB2 = 108~MHz, USART zaman tabanı APB2'den alınır ve
115200~bps baud hatasız üretilir.

\section{Güç ve Topraklama}

\begin{itemize}
    \item STM32 NUCLEO kartı USB üzerinden (ST-Link konektörü) beslenir.
    \item ESP32-CAM \textbf{ayrı bir USB kaynağından} (kendi dock'undan
          veya bağımsız bir USB şarjdan) beslenir. STM32 kartının 5~V
          pininden besleme \emph{brown-out reset}'ine yol açar; kamera +
          Wi-Fi + BLE eşzamanlı çalışırken pik akım yaklaşık 500~mA'dır.
    \item Tek bir kablo bağlantısı yeterli midir? Hayır: STM32 ile ESP-CAM
          arasında en az üç hat olmalıdır --- USART6 TX, USART6 RX ve
          ortak GND. GND referansı paylaşılmadan UART sinyalleri rastgele
          görünür.
    \item ESP32-CAM tarafında flash sırasında IO0 GND'ye çekilir
          (dock üstündeki düğme aynı işi yapar); sonrasında serbest
          bırakılır.
\end{itemize}

\section{Bring-up Sırası}

Bölüm~\ref{tab:bringup} sergiye gitmeden önce takip edilecek bring-up sırasını
özetler. Her satırın ``başarı kriteri'' sütunu sağlandıktan sonra bir
sonraki satıra geçilir.

\begin{table}[h!]
\centering
\caption{Sergi öncesi donanım bring-up sırası ve başarı kriterleri.}
\label{tab:bringup}
\small
\begin{tabularx}{\linewidth}{l X X}
\toprule
\textbf{Adım} & \textbf{İşlem} & \textbf{Başarı Kriteri}\\
\midrule
1 & STM32'yi CubeIDE'den flash et, USART3 VCP'yi Mac terminal'inde aç (115200) & Boot mesajı + LD2 1~Hz heartbeat + B1 USER basınca SCANNING satırı\\
2 & STM32 USART6 TX'i (D1, PG14) bir USB-TTL ile dinle (115200) & B1 USER basıldığında \texttt{CAPTURE\textbackslash n} ve \texttt{SCANNING\textbackslash n} satırları görünür\\
3 & ESP32-CAM'i dock üzerinde flash et (PlatformIO), kendi dock USB'sinden besle & Seri çıktıda Wi-Fi STA IP'si + ``BLE advertising as TaxiGuard'' mesajı\\
4 & ESP32-CAM'i \texttt{eval/far\_frr.py} ile bağımsız test et & EC2 \texttt{/health} 200, sahte burst için MATCH/NOMATCH cevabı\\
5 & STM32 \texttt{D0/D1/GND} pinlerini ESP32-CAM \texttt{IO13/IO14/GND} pinlerine bağla, iki kart farklı USB güçlerinden çalışsın & B1 USER $\rightarrow$ STM \texttt{CAPTURE} $\rightarrow$ ESP burst $\rightarrow$ STM \texttt{RESULT:...} parse\\
6 & iPhone'da TaksiGuvenlik uygulamasını Xcode ile build et, çalıştır & İlk açılışta Login ekranı; kayıt sonrası \texttt{TaxiGuard}'a otomatik bağlantı\\
7 & B1 USER bas, iPhone'da Scan sekmesini izle & Sırasıyla SCANNING $\rightarrow$ MATCH (veya NOMATCH) durum kartı, MATCH'te \texttt{tel://}<numara> dialer'ı açılır\\
8 & Test seti ile FAR/FRR ölçümü yap (Bölüm~\ref{ch:sonuclar}) & ROC eğrisi + EER raporu üretildi\\
\bottomrule
\end{tabularx}
\end{table}
